% Generated draft: requires mathematical and source review.
\documentclass[11pt]{article}
\usepackage[margin=25mm]{geometry}
\usepackage{fontspec}
\setmainfont{Latin Modern Roman}
\usepackage{amsmath,amssymb,mathtools,unicode-math}
\setmathfont{Latin Modern Math}
\usepackage{enumitem,needspace,xurl,hyperref}
\setlength{\parindent}{0pt}
\setlength{\parskip}{5pt}
\setlength{\emergencystretch}{3em}
\newcommand{\N}{\mathbb{N}}
\newcommand{\Z}{\mathbb{Z}}
\newcommand{\Q}{\mathbb{Q}}
\newcommand{\R}{\mathbb{R}}
\newcommand{\C}{\mathbb{C}}
\newcommand{\F}{\mathbb{F}}
\newcommand{\A}{\mathbb{A}}
\newcommand{\PP}{\mathbb{P}}
\newcommand{\E}{\mathbb{E}}
\newcommand{\eps}{\varepsilon}
\newcommand{\dd}{\,\mathrm{d}}
\begin{document}
\section*{UnsolvedMath ID: 701288}
\textbf{Source problem code:} catalogue entry
\subsection{Definitions and mathematical statement}
A group \(G\) is torsion-free if, for every \(g\in G\) and every positive integer \(m\), the equality \(g^m=1\) implies \(g=1\). A subgroup \(H\leq G\) is subnormal of defect at most \(n\) if there exist an integer \(d\) with \(0\leq d\leq n\) and subgroups forming a chain
\[
H=H_0\trianglelefteq H_1\trianglelefteq\cdots\trianglelefteq H_d=G,
\]
where each subgroup is normal in the next. Define the lower central series by \(\gamma_1(G)=G\) and \(\gamma_{i+1}(G)=[\gamma_i(G),G]\), where \([A,B]\) is the subgroup generated by the commutators \(a^{-1}b^{-1}ab\) with \(a\in A\) and \(b\in B\). The condition that \(G\) is nilpotent of class at most \(n\) means \(\gamma_{n+1}(G)=\{1\}\).

For every positive integer \(n\), is it true that every torsion-free group \(G\) for which every subgroup \(H\leq G\) is subnormal of defect at most \(n\) satisfies
\[
\gamma_{n+1}(G)=\{1\}?
\]
The same bound \(n\) is required for all subgroups. No finiteness or finite-generation hypothesis is imposed. The supplied source states that the implication is known to hold for \(n<5\); this qualification is reproduced as a source assertion, without independent verification.
\subsection{Short English statement}
Must a torsion-free group whose every subgroup is subnormal with a common defect bound \(n\) be nilpotent of class at most \(n\)? The supplied record labels the problem open and states that the answer is affirmative for \(n<5\).
\subsection{Sources}
\begin{itemize}
\item Local UnsolvedMath registry: permanent ID \(701288\), exact source code \texttt{catalogue entry}, title \emph{Kourovka Notebook Problem 16.97}. The supplied record labels the problem open.
\item \emph{The Kourovka Notebook (Unsolved Problems in Group Theory)}, original problem 16.97, proposed by G. Traustason in 2006; 16th issue (2006). The supplied background also identifies the 21st issue edition (2026), arXiv:1401.0300v48. The supplied collection description names E. I. Khukhro and V. D. Mazurov as editors.
\end{itemize}
\end{document}
